% =============================================================
%  chapter5_abstract.tex
%  Chapter abstract for Chapter 5 (Draaikraal field study). Extracted
%  verbatim from the original field-report abstract in
%  Chapter5/Workspace/Draaikraal 2023 BJT.tex (the raw pandoc export of
%  Draaikraal Report.docx) -- text unchanged, only the surrounding
%  title/author front matter (handled by chapter5.tex's own \chapter{}
%  and the thesis title page) and the "Abstract" heading itself
%  (handled by the \begin{quote} wrapper in chapter5.tex) were removed.
% =============================================================

The paleoproterozoic Rustenburg Layered Suite (RLS) of the Bushveld
Magmatic Province in northeast South Africa is the largest layered
igneous intrusion in the world, and contains the world's largest source
of platinum group elements as well as important reserves of chromium,
vanadium, titanium, and iron. It is exceptionally well-preserved, and
has been the subject of intense study and debate for generations.
Studies focusing on the intrusion's margin have been especially fruitful
in informing geologists' understanding of igneous petrogenesis and
intrusion mechanisms, modes, and tempos of layered mafic intrusions and
plutonic rocks more generally. The phaneritic RLS is stratigraphically
divided into four zones based on geochemistry and modal mineralogy: the
Lower, Critical, Main, and Upper Zones. These zones are often thought to
have been constructed at differing times and from different magmatic
sources and processes. Thus, the intersection of these stratigraphically
defined zones at the intrusion's margin are particularly interesting
locations for detailed field work. This field study focuses on several
unusual but understudied structural and geomorphic features occurring at
the economically significant intersection of the Main and Critical Zones
with the intrusion's margin in the well-exposed Eastern Limb near the
village of Draaikraal, Mpumalanga, South Africa: the Kliprivier synform,
the Draaikraal Lineament, and the unexpected recessive topography of the
Main Zone in the agricultural valley surrounding the village of
Draaikraal.

The Kliprivier Synform represents a series of modestly rotated plutonic
rocks with conspicuously absent modal layering, and country rocks that
include Quartzite, calc-silicate hornfels, and fine-grained Marginal
Zone gabbros. The country rocks appear to be displaced \textasciitilde{}
1 km into the plutonic RLS near the mapped intersection of the Critical
and Main Zones. The Draaikraal Lineament (named here) is a linear
feature of recessive topography stretching from the country rocks beyond
the intrusion's margin to the middle Main Zone. It overlaps with the
Kliprivier Synform and separates typical, resistant Main Zone
gabbronorites to the north from the recessive Main Zone gabbronorites to
the south in the agricultural valley around Draaikraal. Using structural
measurements and correlation of mottled anorthosite marker beds of the
lower Main Zone, I find that the mottled anorthosite marker beds are
concordant with the floor rocks at the intrusions margin and are not
displaced across the Draaikraal Lineament. I argue that the Kliprivier
Synform is a synmagmatic monocline that formed due to lopolithic
deformation of the country rock roof by the intrusion of Critical Zone
sills before the formation of the UG2 chromitite and overlying Merensky
Reef and Main Zones. The orientation of the Draaikraal Lineament is very
similar to the surface-projection of the axial plane of the Kliprivier
Monocline, which likely continues at depth. The relative recessive
topography of the Draaikraal Lineament and the Main Zone to the south is
argued to be the result of higher degrees of contamination by
hydrothermal fluids than the more resistant Main Zone gabbronorites of
the surrounding area. The rocks at the Draaikraal Lineament are
envisaged to be fluxed with upwardly migrating fluids that were trapped
by the Critical Zone roof, escaping by exploiting the structural
weakness of the Kliprivier Monocline. The agricultural valley to the
south of the lineament is simply envisaged to be relatively contaminated
by melting and devolatilizing sediments at the proximal floor, whereas
the resistant Main Zone to the north was relatively buffered by
interaction with such fluids by the thick Critical Zone rocks that
underlie it there.
